\section{Combinations with repetition}

The last basic scheme combines two ideas: order does not matter, but repetition is allowed. We select a group of objects, and the same kind of object may be selected more than once.

This scheme is less obvious than the previous ones, because the objects we are counting are not naturally written as ordered strings. A selection such as
\[
\{a,a,c\}
\]
contains two copies of \(a\), no copies of \(b\), and one copy of \(c\). The order in which we write these symbols is irrelevant.

\begin{definitionbox}
Let \(A\) be a finite set with \(n\) distinct elements. A combination with repetition of size \(k\) is a selection of \(k\) elements from \(A\), where elements may repeat and order does not matter.
\end{definitionbox}

For example, if
\[
A=\{a,b,c\}
\]
and we select \(2\) elements with repetition allowed, then the possible selections are
\[
\{a,a\},\quad \{a,b\},\quad \{a,c\},\quad \{b,b\},\quad \{b,c\},\quad \{c,c\}.
\]
The selections \(\{a,b\}\) and \(\{b,a\}\) are the same, because order does not matter.

\subsection*{Counting by multiplicities}

A useful way to count such selections is to record only how many times each element is chosen. If the elements are
\[
a_1,a_2,\ldots,a_n,
\]
then a selection of size \(k\) corresponds to a solution of
\[
x_1+x_2+\cdots+x_n=k,
\]
where each \(x_i\) is a non-negative integer. The number \(x_i\) tells how many times \(a_i\) is selected.

For example, if \(A=\{a,b,c\}\), then the selection \(\{a,a,c\}\) corresponds to
\[
x_a=2,\qquad x_b=0,\qquad x_c=1.
\]

\begin{theorembox}
The number of combinations with repetition of size \(k\) chosen from \(n\) distinct elements is
\[
\binom{n+k-1}{k}=\binom{n+k-1}{n-1}.
\]
\end{theorembox}

One way to understand the formula is the stars and bars method. Imagine \(k\) identical stars and \(n-1\) bars separating \(n\) boxes. The stars in the first box count how many times the first element is chosen, the stars in the second box count how many times the second element is chosen, and so on.

For example, with \(n=3\) and \(k=5\), the pattern
\[
\star\star\,|\,\star\,|\,\star\star
\]
represents
\[
x_1=2,
\qquad
x_2=1,
\qquad
x_3=2.
\]
Altogether there are \(k+n-1\) symbols: \(k\) stars and \(n-1\) bars. Choosing the positions of the stars determines the bars, and choosing the positions of the bars determines the stars. Therefore the number of patterns is
\[
\binom{k+n-1}{k}=\binom{k+n-1}{n-1}.
\]

\begin{examplebox}
A bag is to contain \(4\) candies chosen from \(3\) flavours: vanilla, chocolate, and strawberry. Repetition is allowed, and order does not matter. The number of possible bags is
\[
\binom{3+4-1}{4}=\binom{6}{4}=15.
\]
\end{examplebox}

\begin{examplebox}
How many non-negative integer solutions are there to
\[
x_1+x_2+x_3+x_4=7?
\]
This is the same as distributing \(7\) identical stars among \(4\) boxes. Hence the number of solutions is
\[
\binom{7+4-1}{7}=\binom{10}{7}=120.
\]
\end{examplebox}

\begin{remarkbox}
Use combinations with repetition when order does not matter and each type of object may be selected more than once. Typical words are: choose with repetition, distribute identical objects, non-negative integer solutions, multiset.
\end{remarkbox}
